\section{Unbiased estimation after failure and the Gaussian upper bound}\label{sec:measurement}
We use the independence of signs and mapping established in Section~\ref{sec:design} to analyze $\widehat d$ after decoding failure and construct an upper bound for $d$. Cancellation between random signs of different keys makes $\widehat d$ exactly unbiased for every fixed mapping (Proposition~\ref{prop:moments}). It therefore remains unbiased even after arbitrarily rare failures. Controlling the conditional miss rate at the target risk $u$ also requires the distribution of post-failure fluctuations of $\widehat d$ around $d$. When $d\gg\sqrt M$, the standardized estimation error converges to the standard Gaussian distribution (Theorem~\ref{thm:clt}), yielding the Gaussian rule (Definition~\ref{def:gauss}).

\subsection{Conditional unbiasedness and conditional variance given the mapping}
We expand the sum of squares and the total sum in~\eqref{eq:dhat} by key. Each differing key occupies $k$ cells and satisfies $\varepsilon_i^2=1$, so the diagonal contributions of the two terms sum to exactly $d$. The estimation error consists of the cross terms between different keys:
\begin{equation}
\widehat d-d=\frac{2}{k(1-k/M)}\sum_{i<j}\varepsilon_i\varepsilon_j\Bigl(H_{ij}-\frac{k^2}{M}\Bigr).\label{eq:rs-expand}
\end{equation}
The right-hand side is a quadratic form in the random signs. Each coefficient $H_{ij}-k^2/M$ measures the overlap relative to its unconditional mean and depends entirely on the mapping $\mathcal G$. For the standard IBLT, the same expansion replaces the sign products with fixed values $s_is_j$. Selection of overlaps by failure then translates directly into bias (Proposition~\ref{prop:projection}). With random signs, each $\varepsilon_i\varepsilon_j$ takes $\pm1$ with equal probability even on a fixed mapping where failure occurs (Proposition~\ref{prop:composition}). The cross terms therefore have mean zero on every mapping. Intuitively, failure selects mappings with large overlap and amplifies the cross-term coefficients, while each term retains equally likely signs. The error remains centered at zero, and failure changes only its scale. Proposition~\ref{prop:moments} expresses this intuition through the conditional mean and variance given the mapping.

\begin{proposition}[Conditional unbiasedness and conditional moments given the mapping]\label{prop:moments}
For every fixed mapping $\mathcal G$,
\begin{equation}
\E[\widehat d\mid\mathcal G]=d,\qquad
\Var(\widehat d\mid\mathcal G)=\frac{4}{k^2(1-k/M)^2}\sum_{i<j}\Bigl(H_{ij}-\frac{k^2}{M}\Bigr)^2\le2d(d-1).\label{eq:condmoments}
\end{equation}
Thus, for any mapping event $\mathcal S$ of positive probability, $\E[\widehat d\mid\mathcal S]=d$. In particular, $\E[\widehat d\mid F]=d$ holds for every finite $M$ and every sign composition, even for arbitrarily rare failures. Averaging over the random mapping $\mathcal G$ gives $\E\widehat d=d$ and $\Var\widehat d=s(d)^2$.
\end{proposition}

\begin{proof}[Proof sketch]
Given $\mathcal G$, $(\varepsilon_i)$ are iid uniform signs. Taking expectations term by term in~\eqref{eq:rs-expand} and using $\E[\varepsilon_i\varepsilon_j\mid\mathcal G]=0$ proves conditional unbiasedness. The sign products for distinct unordered pairs $(i,j)$ are pairwise uncorrelated, yielding the conditional variance in~\eqref{eq:condmoments}. The Cauchy--Schwarz inequality gives the upper bound $2d(d-1)$. The full proof is in Appendix~\ref{app:proof-moments}.
\end{proof}

For a fixed mapping, $\widehat d$ is a Rademacher quadratic form with the Gram matrix of the centered incidence vectors. Its conditional mean equals the trace, and its conditional variance comes only from the off-diagonal terms. These are the same identities used by the Hutchinson trace estimator\cite{hutchinson1990trace}. Proposition~\ref{prop:moments} expresses them in terms of overlap coefficients, establishing conditional unbiasedness for any mapping event.

Figure~\ref{fig:composition} illustrates this unbiasedness numerically. The post-failure bias of the standard encoding changes direction with $k$, while the mean for the random-sign encoding stays near zero. Section~\ref{sec:exp-decode} and Appendix~\ref{app:dist} compare conditional moments across additional configurations.

\subsection{\texorpdfstring{The Gaussian limit after failure when $d\gg\sqrt M$}{The Gaussian limit after failure when d >> sqrt(M)}}
Proposition~\ref{prop:moments} establishes the center of the estimator after failure. Constructing an upper bound also requires the shape of the error distribution around this center. Failure affects $\widehat d$ only through the mapping. By~\eqref{eq:condmoments}, the sum of squared overlaps $\sum_{i<j}(H_{ij}-k^2/M)^2$ determines the conditional error scale. Failure selects mappings with different values of this sum and may therefore change the error distribution. We standardize the error by the unconditional standard deviation $s(d)$:
\begin{equation}
X=\frac{\widehat d-d}{s(d)},\label{eq:X}
\end{equation}
and study the conditional distribution of $X$ given $F$.

We first fix a mapping. By~\eqref{eq:rs-expand}, the error is a weighted sum of sign products $\varepsilon_i\varepsilon_j$. The main contributions come from key pairs sharing a cell: the intersecting edge pairs with $H_{ij}\ge1$ in $\mathcal G$. Their expected number is about $k^2d^2/(2M)$. Sign products for distinct key pairs are pairwise uncorrelated, and no small set of terms dominates the sum. The number of intersecting edge pairs grows with $d^2/M$ and tends to infinity when $d\gg\sqrt M$. The error then sums many $\pm1$ perturbations, and its distribution approaches a Gaussian. When $d$ is of the same order as $\sqrt M$, only finitely many intersecting edge pairs contribute. The error consists of a few $\pm1$ terms, and its distribution remains discrete. We distinguish these regimes by the scale parameter
\begin{equation}
c=\frac{d}{\sqrt M}.\label{eq:c}
\end{equation}
The expected number of intersecting edge pairs is about $k^2c^2/2$, so $c$ plays the role of the collision scale in the birthday paradox\cite{arratia1989chen,dw}. It and the load $d/M$ are independent parameters.

Conditioning on failure preserves the sign structure. On any mapping where failure occurs, $(\varepsilon_i)$ remain iid uniform signs (Proposition~\ref{prop:composition}). The Gaussian shape persists as long as the mapping still contains many intersecting edge pairs. This also holds for rare failures. Under a subcritical load, a typical failure is caused by a duplicate edge pair. This raises a single $H_{ij}$ to $k$ and leaves the overall structure of intersecting edge pairs almost unchanged.

\begin{theorem}[Gaussian limit after failure]\label{thm:clt}
Fix $k\ge3$ and let $d/\sqrt M\to\infty$, allowing any limiting behavior of $d/M$. Then
\begin{equation}
\sup_x\bigl|\Pr(X\le x\mid F)-\Phi(x)\bigr|\longrightarrow0.\label{eq:clt}
\end{equation}
This includes the critical window, rare failures under a fixed subcritical load, and arbitrary overload.
\end{theorem}

\begin{proof}[Proof sketch]
The proof has two steps. First, we fix the mapping and take randomness only over the signs. Suppose the number of intersecting edge pairs is of the same order as $M$ and each key participates in overlaps of size at most $O(\log M)$. The higher cumulants of the standardized quadratic form~\eqref{eq:rs-expand} then tend to zero, and the conditional distribution of $X$ converges to $\mathcal N(0,1)$ uniformly over such mappings. Second, we show that, conditional on $F$, the mapping satisfies these conditions with probability tending to $1$. We consider three cases according to the limit of $d/M$. When $d/M$ has a positive finite limit, $\Pr(F)\ge\binom Mk^{-1}$. The probability that a mapping violates the conditions is superpolynomially small and remains negligible after conditioning. When $d=o(M)$, Proposition~\ref{prop:plant} shows that the mapping distribution given failure is asymptotically identical to the planted distribution in TV distance. The Gaussian limit holds for the planted distribution. When $d>M$, failure occurs with certainty, so the conditional and unconditional distributions coincide. The full proof is in Appendix~\ref{app:clt}.
\end{proof}

Figure~\ref{fig:encoding} compares the distributions of $X$ for failure samples and all samples in a configuration with $c\approx104$. For the standard encoding, we compute $X$ in the same way from its counts $C_v=\sum_{i:v\in S_i}s_i$. With random-sign encoding, the two distributions almost coincide. With standard encoding, the failure distribution is shifted as a whole.

\subsection{The Gaussian rule and its conditional miss rate}\label{sec:gauss}
By Theorem~\ref{thm:clt}, $\widehat d$ is approximately Gaussian after failure, with mean $d$ and standard deviation $s(d)$. If $d$ were known, we could set the margin using the fact that $\widehat d$ falls below $d+s(d)\Phi^{-1}(u)$ with probability about $u$. However, the scale $s(d)$ (\eqref{eq:s}) depends on the unknown $d$. Direct substitution of $s(\widehat d)$ would make the margin fluctuate with the observation. A low $\widehat d$ gives a small $s(\widehat d)$, reducing the margin precisely for samples that need a larger one. The risk analysis would then need to account for this additional error.

Test inversion is a classical method for constructing confidence sets\cite{neyman1937}. We reverse the reasoning and test candidate values individually. For each candidate $n$, we assume the true value is $n$. The scale $s(n)$ is then known, and $\widehat d$ falls below $n+s(n)\Phi^{-1}(u)$ with probability about $u$. An observation below this lower bound indicates that $n$ is too large, so we exclude $n$. The largest remaining candidate gives the upper bound $U$. The risk is controlled by a single event: $U<d$ implies that the true value $d$ is excluded. Under the true value $d$, this exclusion has probability about $u$. Definition~\ref{def:gauss} follows this construction, replacing $\Phi^{-1}(u)$ with a more accurate quantile for finite $M$.

\begin{definition}[Gaussian rule]\label{def:gauss}
Given an observation $\widehat d$ and a target risk $u\in(0,1/2)$, we test each candidate integer $n\ge2$ as the true value, for which the scale $s(n)$ is known. Let $\chi^2_{M-1,u}$ denote the lower $u$ quantile of the chi-square distribution with $M-1$ degrees of freedom. Set the acceptance lower bound to
\begin{equation}
t(n)=n+s(n)\,\frac{\chi^2_{M-1,u}-(M-1)}{\sqrt{2(M-1)}}.\label{eq:gauss}
\end{equation}
We accept candidate $n$ if $\widehat d\ge t(n)$ and reject it otherwise. We set $U$ to the largest accepted candidate, returning $+\infty$ if none is accepted.
\end{definition}

For fixed $M,k$ as $d\to\infty$, $\widehat d/d$ converges in distribution to $\chi^2(M-1)/(M-1)$\cite{wu2026selfsizing}, which has mean 1. After centering the chi-square variable at its mean and scaling it to unit variance, its lower $u$ quantile is $(\chi^2_{M-1,u}-(M-1))/\sqrt{2(M-1)}$. Equation~\eqref{eq:gauss} uses this quantile in place of $\Phi^{-1}(u)$ to capture skewness more accurately for finite $M$. The two quantiles agree as $M\to\infty$. In the experimental range $M\ge64$, $u\ge0.01$, $t(n)$ is strictly increasing for $n\ge2$ (Appendix~\ref{app:gauss-mono}), so we can find $U$ by monotone binary search. Outside this range, we check candidates individually as specified in the definition. The rule depends only on the observation $\widehat d$ and the known parameters $(M,k,u)$.

\begin{corollary}[Conditional miss rate of the Gaussian rule]\label{cor:gauss}
Under the conditions of Theorem~\ref{thm:clt}, for any fixed $u\in(0,1/2)$, $\limsup\Pr(U<d\mid F)\le u$.
\end{corollary}

\begin{proof}[Proof sketch]
We take $n=d$. Theorem~\ref{thm:clt} and $(t(d)-d)/s(d)\to\Phi^{-1}(u)$ give a rejection probability of $u+o(1)$ for the true candidate. The conclusion follows from $\{U<d\}\subseteq\{\widehat d<t(d)\}$. The full proof is in Appendix~\ref{app:clt}.
\end{proof}

By Corollary~\ref{cor:gauss}, the Gaussian rule has conditional miss rate asymptotically at most $u$ throughout the range $d\gg\sqrt M$ of Theorem~\ref{thm:clt}. This range includes the critical window and the overload region. Figure~\ref{fig:coverage} shows that the Gaussian rule's miss/u approaches the target level as $c$ grows. When $d$ is of the same order as $\sqrt M$ or smaller, the Gaussian approximation no longer holds, and the actual conditional miss rate exceeds $u$. Section~\ref{sec:inversion} explains this behavior and gives the Skellam rule, whose conditional miss rate is asymptotically at most $u$ uniformly over all $d\ge2$.
